\documentclass[preprint,12pt]{elsarticle}

\usepackage{booktabs}
\usepackage{url}
\usepackage[hidelinks]{hyperref}
\usepackage{xcolor}

\newcommand{\pendiente}[1]{\textcolor{red}{[\textbf{pendiente:} #1]}}

\journal{Computer Standards \& Interfaces}

\begin{document}

\begin{frontmatter}

\title{Conformance vectors written from the specification find defects that
       implementation-derived test suites do not: evidence from an HTTP
       server}

\author[una]{Richar Andre Vilca-Solorzano\corref{cor}}
\ead{75521963@est.unap.edu.pe}
\ead[orcid]{0009-0003-2385-5263}
\cortext[cor]{Corresponding author.}

\author[una]{Fred Torres-Cruz}
\ead[orcid]{0000-0003-0834-6834}

\author[una]{Ramiro Pedro Laura-Murillo}
\ead[orcid]{0000-0003-1837-4871}

\affiliation[una]{organization={Universidad Nacional del Altiplano,
                                Facultad de Ingeniería Estadística e Informática},
                  city={Puno}, country={Peru}}

\begin{abstract}
\pendiente{escribir al final}
\end{abstract}

\begin{keyword}
HTTP \sep protocol conformance \sep request smuggling \sep software testing \sep
RFC 9112 \sep RFC 9113 \sep web server security
\end{keyword}

\end{frontmatter}

\section{Introduction}
\label{sec:introduction}

A test suite written alongside an implementation tends to assert what the
implementation was built to do. It is a poor instrument for finding what the
specification requires and the implementation never considered.

This paper reports a controlled instance of that gap. We took an HTTP server
that passed its own suite of 434 checks and subjected it to two independent
conformance instruments: a set of vectors we wrote directly from the text of
RFC~9112 \cite{rfc9112} and RFC~9110 \cite{rfc9110}, and h2spec
\cite{h2spec}, an existing suite for HTTP/2 \cite{rfc9113}. We report what each
instrument found, what an external security audit found, and --- the part we
believe generalises --- what distinguishes the defects that the
implementation-derived suite could never have found.

\paragraph{Contributions}
\begin{enumerate}
  \item A reusable set of 23 conformance vectors for HTTP/1.1 servers covering
        13 clauses of RFC~9112 and RFC~9110, sent as raw bytes and each citing
        the clause it enforces. No free conformance suite for HTTP/1.1 servers
        exists, unlike for HTTP/2.
  \item The four violations the vectors exposed in an implementation that
        already passed its own tests (Section~\ref{sec:violations}).
  \item Two defects of the request-smuggling family whose common shape is an
        \emph{unexercised exit path}, not an unexercised line
        (Section~\ref{sec:exitpaths}), one of which was intermittent and
        therefore invisible to a pass/fail criterion applied once.
  \item A characterisation of which defect classes each instrument detects, and
        the cost of each (Section~\ref{sec:discussion}).
\end{enumerate}

\section{Background and related work}
\label{sec:background}

\pendiente{revisión de literatura: conformidad de protocolos, request
smuggling, fuzzing de servidores HTTP, suites de conformidad existentes}

\section{Subject system and instruments}
\label{sec:subject}

\subsection{Subject}

Cero is a Java web framework with its own HTTP/1.1 and HTTP/2 server and no
runtime dependencies beyond the JDK. The HTTP module comprises 43 classes and
6\,752 lines, exercised by 434 checks of the project's own suite. The complete
system is public.

\subsection{Instruments}

\begin{table}[ht]
\centering
\caption{The four instruments applied, and what each one is.}
\label{tab:instruments}
\begin{tabular}{lll}
\toprule
Instrument & Origin & Size \\
\midrule
Project test suite            & written with the implementation & 434 checks \\
Conformance vectors (HTTP/1.1) & written from the RFC text      & 23 vectors, 13 clauses \\
h2spec (HTTP/2)               & third party \cite{h2spec}       & 146 tests \\
External security audit       & independent reviewer            & 11 findings \\
\bottomrule
\end{tabular}
\end{table}

The distinction that matters is the second column. Only the first instrument
had access to the implementation while being written.

\section{Violations found by specification-derived vectors}
\label{sec:violations}

The 23 vectors exposed four violations in a server that passed its own suite:

\begin{enumerate}
  \item \textbf{\emph{absolute-form} request targets rejected} (RFC~9112 §3.2.2).
        A server must accept \texttt{GET http://host/path HTTP/1.1}; ours
        returned 400.
  \item \textbf{\emph{asterisk-form} rejected} (RFC~9112 §3.2.4).
        \texttt{OPTIONS * HTTP/1.1} was not recognised.
  \item \textbf{\texttt{chunked} accepted outside the final position} of
        \texttt{Transfer-Encoding} (RFC~9112 §6.1). Accepting it is a known
        precondition for request smuggling when an intermediary disagrees with
        the origin about where a message ends.
  \item \textbf{Null bytes admitted in header values} (RFC~9110 §5.5).
\end{enumerate}

None of the four is exotic. All four survived a suite written by the same
authors who wrote the parser, because that suite asserted the behaviour the
parser was designed for.

\section{Unexercised exit paths}
\label{sec:exitpaths}

Two further defects share a shape that we argue is the more interesting
finding, because neither is reachable by adding assertions to the existing
suite.

\subsection{The session cookie over HTTP/2}

Two sibling implementations of the same response interface existed, one per
protocol version. Only the HTTP/1.1 one asked for the pending session cookie.
Over HTTP/2, therefore, no client could establish a session; and with no
session there is no CSRF token, so every write returned 403 attributing the
failure to the wrong cause.

All 432 checks of the module passed throughout. Every one of them entered over
HTTP/1.1.

\subsection{Body validation after the response}

Completing a response closes the local half of an HTTP/2 stream, not the stream
(RFC~9113 §5.1) \cite{rfc9113}. The implementation removed the stream from its
table as soon as the request handler returned. A \texttt{DATA} frame arriving
afterwards therefore found no stream, and the check that the received body
matches the declared \texttt{content-length} --- the check that prevents the
smuggling precondition --- was skipped.

The defect was \emph{intermittent by construction}: it manifested only when the
response won the race against the request body. h2spec failed approximately one
run in five; the project suite passed always. A conformance tool run once in
continuous integration is a sampling instrument, not a decision procedure, when
the defect is a race.

\subsection{Why line coverage cannot see either}

In both cases every line was executed by some test. What was never traversed was
a \emph{path out of the system}: a response emitted over the second protocol, a
frame arriving after the handler returned. We therefore propose counting exit
paths --- protocol version $\times$ response kind $\times$ ordering --- as the
coverage criterion for this class of component.

\section{External audit}
\label{sec:audit}

The framework's first external consumer was a university portal. Its developer,
who had not written the framework, reviewed the source and produced eleven
findings, two of them exploitable from outside without credentials. Six are
security defects and are listed in Table~\ref{tab:audit}; the remaining five
were API and functional gaps.

\begin{table}[ht]
\centering
\caption{Security findings from the external review. None was reachable by the
         project's own suite at the time.}
\label{tab:audit}
\begin{tabular}{p{0.26\linewidth}p{0.62\linewidth}}
\toprule
Finding & Mechanism \\
\midrule
Session fixation &
  The identifier held before authentication remained valid afterwards; the CSRF
  token travelled with it. \\
Memory exhaustion via static files &
  The classpath resource cache stored \emph{absences}, keyed by the requested
  path. Every 404 left two permanent entries in an unbounded map, so an
  unauthenticated loop filled the heap. \\
Rate limiter bypass &
  The path formed part of the bucket key, so spreading the load across invented
  URLs multiplied the quota while growing the map. \\
Symbolic links followed &
  \texttt{readAttributes} followed them although the adjacent line declared
  otherwise. \\
No absolute session lifetime &
  Only inactivity was counted, so a session touched occasionally never expired. \\
CSRF exemption by bare prefix &
  Exempting \texttt{/api/public} also exempted
  \texttt{/api/publicSECRET}. \\
\bottomrule
\end{tabular}
\end{table}

Two of these --- the unbounded cache and the rate-limiter key --- are resource
exhaustion reachable without credentials. Both arise from the same modelling
error: a cache key chosen from attacker-controlled input with no bound on
cardinality.

Two additional defects appeared only after deployment behind a reverse proxy: a
session cookie issued without \texttt{Secure}, because the application did not
know TLS was terminated upstream, and a deployment script that erased the
certbot HTTPS block when rewriting the virtual host. Neither is reachable by any
instrument applied to the program alone.

\section{Results}
\label{sec:results}

\pendiente{tabla resumen: defectos por instrumento y por clase, con el coste de
cada instrumento en horas}

\section{Discussion}
\label{sec:discussion}

\pendiente{qué clase de defecto detecta cada instrumento; coste-beneficio;
recomendación para implementadores de servidores}

\section{Threats to validity}
\label{sec:threats}

\paragraph{Single subject}
The findings come from one implementation. We report the shape of each defect
rather than a frequency, and make no claim about how often these defects occur
in HTTP servers generally.

\paragraph{The vectors are ours}
The 23 vectors were written by the same authors as the implementation, which
weakens their independence. We mitigated this by writing them from the RFC text
without reference to the code, and by cross-checking HTTP/2 against a suite we
did not write \cite{h2spec}. The audit is the only fully independent instrument.

\paragraph{Intermittent defects and sampling}
Our estimate of one failure in five runs for the race described in
Section~\ref{sec:exitpaths} is based on \pendiente{número de corridas} runs and
is reported as an observation, not a measured rate.

\section{Conclusion}
\label{sec:conclusion}

\pendiente{cerrar}

\section*{CRediT authorship contribution statement}

\textbf{Richar Andre Vilca-Solorzano:} Conceptualization, Methodology,
Software, Validation, Investigation, Writing --- original draft.
\textbf{Fred Torres-Cruz:} Conceptualization, Methodology, Writing --- review
\& editing, Supervision. \textbf{Ramiro Pedro Laura-Murillo:}
Conceptualization, Writing --- review \& editing, Supervision.

\section*{Declaration of competing interest}

The authors declare no competing financial interests or personal relationships
that could have appeared to influence the work reported in this paper.

\section*{Data availability}

The implementation, its test suite, the 23 conformance vectors and the scripts
that reproduce every result are public at
\url{https://github.com/Andre031222/Cero}.

\bibliographystyle{elsarticle-num}
\bibliography{referencias}

\end{document}
